\begin{derivation}[the click's algebra beyond the law, Section 2.5.2]
\emph{Inputs:} Rule3 at the vacuum's paces and its band $\cos\omega = (\num / 3\den)\sum_a \cos k_a$ (S.1); the readings at a Node, the Wronskian $W = \mathrm{re}_{\mathrm{now}}\,\mathrm{im}_{\mathrm{before}} - \mathrm{im}_{\mathrm{now}}\,\mathrm{re}_{\mathrm{before}}$ of a plane configuration and the Link current $F_{ij} = \num\,(\mathrm{now}_i\,\mathrm{before}_j - \mathrm{before}_i\,\mathrm{now}_j)$ of a real line (Section 3.5); the invariant $2A^2\sin\omega = T$ per quantum and the energy $T\sin\omega$ (Section 4.3); light's speed $c = 1 / \sqrt 3$; and the proposition of Section 2.5.2, a line beyond the law's write under its own name, that the write carries the arriving wave vector $\mathbf k$ as a phase twist of the taker's record, $\delta k_a = k_a\,\mathrm{div}\,M$ per Link along each axis for a record of count $M$, the remainder $k_a\,\mathrm{mod}\,M$ carried at the body. Nothing below is read from a run.
\emph{Steps:} (a) \emph{Lemma A.} For a plane configuration $a_i = A\,e^{i(\mathbf k\cdot\mathbf x_i - \omega t)}$, with $\mathrm{now} = A e^{i\phi}$ and $\mathrm{before} = A e^{i(\phi + \omega)}$ one interval earlier, $W = A^2[\cos\phi\sin(\phi + \omega) - \sin\phi\cos(\phi + \omega)] = A^2\sin\omega$ at every Node, free of $\mathbf k$; for a real line along the axis $a$ with $\phi_j = \phi_i + k_a$, $\mathrm{now}_i\mathrm{before}_j - \mathrm{before}_i\mathrm{now}_j = A^2[\cos\phi\cos(\phi + k_a + \omega) - \cos(\phi + \omega)\cos(\phi + k_a)] = \tfrac12 A^2[\cos(k_a + \omega) - \cos(\omega - k_a)] = -A^2\sin\omega\sin k_a$, so $F_{ij} = -\num\,A^2\sin\omega\sin k_a$, free of $\phi$. A phase common to the two levels at a Node leaves $W$; a phase step across a Link moves $F$. (b) \emph{The norm factor.} The one-Node recurrence $x_{t+1} + x_{t-1} = 2\cos\omega\,x_t$ conserves $Q = x_t^2 + x_{t-1}^2 - 2\cos\omega\,x_t x_{t-1}$: for a real line $B\cos(\phi - \omega t)$, $Q = B^2\sin^2\omega$; for a plane configuration of amplitude $A$, $Q = 2A^2\sin^2\omega$, so $Q / \sin\omega = 2A^2\sin\omega = T$ and $W = A^2\sin\omega = T / 2$ per plane quantum. The identity $2A\cos\phi = Ae^{i\phi} + Ae^{-i\phi}$ is therefore exact in the count, $4A^2\sin^2\omega$ on both sides: a real line of amplitude $2A$ is two plane quanta of opposite sense, neutral, $W = +A^2\sin\omega - A^2\sin\omega = 0$; one real quantum has $B_q^2 = 2A_q^2$, and a real line of amplitude $B$ carries $B^2 / (2A_q^2)$ quanta. (c) \emph{The twist.} Under the proposition the twist $e^{i\,\delta\mathbf k\cdot\mathbf x}$ is a phase common to the two levels at each Node, so by (a) it leaves $W$ at every Node, the charge kept, and moves the Link current; it keeps the count whole only with the write's own scaling of the amplitude, Section 2.5's $A^2 = (N + 1)\,T / (2\sin\Omega)$ taken at the twisted rotation, $A^2 = (N + 1)\,T / (2\sin\omega(k / M))$, the invariant $2A^2\sin\omega = T$ per quantum holding at every $k$ and the energy's difference $T[\sin\omega(k) - \sin\omega_0]$ being the recoil of (e), booked by the write; the integer wave vector, in units of $2\pi / N$ on a box of $N$ Nodes, is conserved across the click modulo $2\pi$ per Link with the remainder carried, while the momentum reading $F \propto \sin k_a$ is conserved only to the order $k^3$ and folds at the zone's edge. The wave vector the twist carries is read at the taker's Ports and nowhere else: by (a) the arriving light's Link current along the axis $a$ is $F_{ij} = -\num\,B^2\sin\omega\sin k_a$ with $B^2\sin\omega = T$ per real quantum (b), so $\sin k_a$ per axis is in the NodeDetector's own reads over its boundary, the integer $k_a$ its nearest whole in the box's units to the order $k^3$ of the reading, and no number beyond the reads enters. (d) \emph{The band's inertia and kinetic scale.} With $\cos\omega_0 = \num / \den = \nu$ and $\sum_a(1 - \cos k_a) = k^2 / 2 + O(k^4)$, $\cos\omega - \cos\omega_0 = -\nu k^2 / 6$, so $\omega - \omega_0 = \nu k^2 / (6\sin\omega_0) = k^2 / (2m^*)$ with $m^* = 3\sin\omega_0 / \nu = 3\tan\omega_0$, and $\omega^2 = \omega_0^2 + c_m^2 k^2 + O(k^4)$ with $c_m^2 = \omega_0 / (3\tan\omega_0) = c^2\,\omega_0 / \tan\omega_0$ (S.12; $c_m / c = 0.867$ at $[2, 3]$). (e) \emph{The recoil.} A body of count $M$ taking the twist $k / M$ per Link gains $M[\omega(k / M) - \omega_0] = k^2 / (2Mm^*) + O(M^{-3})$, vanishing for a large body; after an absorption and an emission at $\mathbf k_{\mathrm{out}}$ the body keeps $\mathbf k_{\mathrm{in}} - \mathbf k_{\mathrm{out}}$ exactly. (f) \emph{Compton.} The click balances the rotations by the beat rule of Section 4.3, the law's own (the born light at $\omega_e - \omega_g$, the energy in the click's unit $T\sin\omega$ following to the order $\omega_i\omega_j$ while no layer keeps a total, Section 2.4), and under the proposition the wave vector; balanced in the rotations the non-relativistic factor below is $\omega_0 / \tan\omega_0$, where a balance of the energies $T\sin\omega$ would give $\omega_0\cos^2\omega_0 / \sin\omega_0 = 0.501$ at $[2, 3]$, the two one at nature's gap; with light's band $\Omega = c|\mathbf K|$ and the matter band of (d) in the Klein-Gordon form, a quantum at rest absorbing $\Omega_{\mathrm{in}}$ and emitting $\Omega_{\mathrm{out}}$ at the angle $\theta$ obeys $(1 - \beta)(\Omega_{\mathrm{in}}^2 + \Omega_{\mathrm{out}}^2) - 2\Omega_{\mathrm{in}}\Omega_{\mathrm{out}}(1 - \beta\cos\theta) + 2\omega_0(\Omega_{\mathrm{in}} - \Omega_{\mathrm{out}}) = 0$ with $\beta = c_m^2 / c^2 = \omega_0 / \tan\omega_0$. At $\beta = 1$, nature's gap, this is $1 / \Omega_{\mathrm{out}} - 1 / \Omega_{\mathrm{in}} = (1 - \cos\theta) / \omega_0$, Compton's formula, $\lambda_{\mathrm{out}} - \lambda_{\mathrm{in}} = (2\pi c / \omega_0)(1 - \cos\theta)$ with $\lambda = 2\pi c / \Omega$, the Compton wavelength $2\pi / (\sqrt 3\,\omega_0)$ Links. At a finite gap the outgoing quantum gains, to first order in $1 - \beta$, $\delta\Omega_{\mathrm{out}} = (1 - \beta)\,c^2 q^2\,\Omega_C / (2\omega_0\Omega_{\mathrm{in}})$, $q$ the recoil wave vector and $\Omega_C$ Compton's value; in the non-relativistic regime $\Omega_{\mathrm{in}} \ll \omega_0$ this is one factor on the whole shift, $\lambda_{\mathrm{out}} - \lambda_{\mathrm{in}} = (2\pi / (m^* c))(1 - \cos\theta) = (2\pi / (\sqrt 3\,\omega_0))\,(\omega_0 / \tan\omega_0)\,(1 - \cos\theta)$, exact at every gap in that regime: the inertia $3\tan\omega_0$ against the rest rotation $\omega_0$, $\omega_0 / \tan\omega_0 = 1 - \omega_0^2 / 3 + O(\omega_0^4) = 1 - \tfrac23(1 - \nu) + O((1 - \nu)^2)$, $0.752$ at $[2, 3]$ and $1$ at nature's gap. (g) \emph{Absorption alone.} A free quantum at rest absorbing $\Omega = cK$ with no emission would need $cK = \sqrt{\omega_0^2 + c_m^2 K^2} - \omega_0$, that is $K(c^2 - c_m^2) = -2c\,\omega_0 < 0$ for $c_m \le c$: forbidden at every gap, the free quantum's one window holding two writes. (h) \emph{The zone's edge.} A single real quantum making a plane pair needs $K^2(c^2 - c_m^2) \ge 4\omega_0^2$: impossible at nature's gap, and at $[2, 3]$ only at $K \ge 2\omega_0 / \sqrt{c^2 - c_m^2} = 5.85$ radians per Link, outside the zone's $|K| \le \pi\sqrt 3 = 5.44$.
(i) \emph{The integer recoilless condition.} The twist is an integer division, so $\delta k_a = k_a\,\mathrm{div}\,M = 0$ whenever $k_a < M$ in the box's units of $2\pi / N$ per Link: a body whose count exceeds the quantum's wave number takes no twist from one click; the remainder $k_a\,\mathrm{mod}\,M = k_a$ accumulates at the body, and one whole unit per Link is written at every $M / k_a$ clicks on average, the first at the $\lceil M / k_a\rceil$-th, the momentum conserved exactly over the accumulation and never per click. The threshold $k_a < M$ compares a wave number in the declared box's units with a count, so it is a property of the computed world and no formula shared with nature; the box-free statement is the recoil energy of (e). (j) \emph{Light's wave vector in a well.} At static paces Rule3 keeps a wave's frequency, and light's local speed in a well of potential level $U$ is $c / n$ with the index $n = e^{2U}$ (S.2), so its wave vector is $k = \omega n / c$ at long wavelength (the paper's Eq. for the lens, the chromatic term $k^2(e^{6U} - e^{2U}) / 24$ beyond it; S.28): under the proposition the recoil per click inside the well is $n$ times the vacuum's at the same frequency, Minkowski's form of the momentum of light in a medium.
\emph{Check:} \texttt{click\_algebra.py} evaluates (a) and (b) in floating point to $10^{-15}$, (d) at $k = 0.05$, (f) against the exact quadratic at $\Omega_{\mathrm{in}} = 0.02$ and $\theta = \pi / 2$ (the ratio $0.757$, the factor $0.752$ and the relativistic remainder), (g) and (h) at $[2, 3]$, and the accumulation of (i) for $k_a = 3$ and $M = 7$ (the kicks at the third, fifth and seventh clicks); one Node and one Link, no run.
\emph{Status:} (a), (b), (d) derived; (c), (e), (f), (g), (h), (i), (j) derived under the proposition of Section 2.5.2 (hypothesis, fitted); lattice, with S.2 and S.28 for (j); no run; (i)'s threshold a property of the declared box and no formula shared with nature. The angular shares of the emitted direction are not derived here.
\end{derivation}
